% chapter/chapter2.tex
% The Unquiet Dead – Echoes, Anchors, and the Hungry Silence
% Rewritten in the voice of Lerris Fenwood, with marginalia from the Gray Wanderer,
% Saikou Ira, and Dusana of the Raven Road.

\epigraph{“The dead do not forget. They do not forgive. But they will negotiate – if you remember to count to eight before you speak.”}{—Saikou Ira, from a case file marked “closed – unsatisfactory”}

The dead are not a single kind of hunger. I learned this in the Wardens’ archive, when a leaf that had been catalogued as “unremarkable” began to weep a thin, dark fluid that smelled of iron and old rain. The cataloguer had left the ninth line blank. The dead do not need violence; they need a witness.

This chapter catalogs the unquiet: from the faint echo of a scream repeated at midnight to the poltergeist that has learned to count your breaths. Saikou Ira taught me that a ghost is not a monster; it is a debt that has forgotten how to be paid. I have tried to honor that distinction here.

\section*{The Spirit Taxonomy (After Saikou Ira)}

Ira’s taxonomy is not a scientific classification; it is a \textbf{pragmatic ledger} of how the dead manifest and how they may be quieted. I have reproduced it with his permission – and with the caveat that no ghost ever read a taxonomy before deciding to haunt.

\subsection*{Class I: The Echo}

\noindent\textbf{Description}: Not truly a spirit – a recording. A psychic imprint left by a traumatic death, repeated endlessly. Echoes do not speak, do not learn, and cannot be negotiated with. They are the scream that plays again at midnight, the footsteps that cross the same corridor at the same hour.

\noindent\textbf{Stats}: Tier I, Cap 1 (non‑combatant). Incorporeal, Repetitive, Trapped.\\
\textbf{Harm}: Cannot be harmed directly. \textbf{Fatigue}: None.

\noindent\textbf{Traits}: 
\emph{Incorporeal} – cannot be touched, cannot touch. \emph{Repetitive} – performs the same sequence of actions every cycle, with no variation. \emph{Trapped} – cannot leave the site of its imprint.

\noindent\textbf{Complications}: Distraction (SB: 1 to create a sensory illusion), Delay (SB: 2 to cause a time loop – the same moment repeats for one exchange).

\noindent\textbf{Resolution}: Witness a truth about the death (DV 3) to banish permanently. A single observer who speaks aloud the truth of the event – who names the dead and describes their end – will often break the loop. I have done this twice. The first time, the Echo thanked me. The second time, it was already gone.

\noindent\textbf{Clock}: \emph{Echo’s Loop [4]} – when filled, the Echo expands its range. A scream that once stayed in the attic will begin to walk the stairs.

\saikou{I once placed a silver mirror facing the site of an Echo. The scream played once into the glass. When I shattered the mirror, the Echo did not return. But the shards held fragments of the memory. I keep them in a lead box.}

\subsection*{Class II: The Anchor}

\noindent\textbf{Description}: A spirit bound to a specific object, location, or person by unresolved emotion – grief, rage, love, duty. Unlike the Echo, an Anchor possesses will and memory. It can speak, bargain, and choose to manifest. It is, in essence, a person who has forgotten how to leave.

\saikou{There is a ghost near a ditch in Thepyrgos... murdered and anchored as he was encircled with salt when he died. His soul could not escape. He wails for relief that will never come; the one who anchored him is long dead. He was a wretch, his tale his known. Some buisiness should remain unfinished.}

\noindent\textbf{Stats}: Tier II‑III, Cap 2–4. Incorporeal, Intelligent, Bound.\\
\textbf{Harm}: 3. \textbf{Fatigue}: Spirit.

\noindent\textbf{Traits}: 
\emph{Manifestation} (DV 3, Fatigue 1) – the Anchor becomes visible and may speak. \emph{Poltergeist Shove} (DV 4, Harm 1) – telekinetic force.

\noindent\textbf{Complications}: Bargain Required – cannot be fought; must negotiate. The Anchor has a price: a confession, a returned object, a spoken name.

\noindent\textbf{Resolution}: Fulfill the Anchor’s unfinished business, or perform a witnessed substitution ritual (DV 4). I once helped an Anchor complete a letter to a daughter who had already died; the ghost read the letter aloud and then – did not vanish, but smiled. That was enough.

\noindent\textbf{Clock}: \emph{Anchor’s Patience [6]} – when filled, the Anchor becomes a Poltergeist (Class IV). The patience of the dead is not infinite; it is merely slow.

\dusana{An Anchor in the Mistlands asked me for a cup of tea. I poured it. She told me the name of the man who had drowned her. I wrote it down. She did not pass; she simply stopped being sad.}

\subsection*{Class IV: The Poltergeist}

\noindent\textbf{Description}: An Anchor that has been corrupted by rage, starvation, or the failure of negotiation. It is no longer a person – it is a wound that has learned to strike back. Poltergeists manifest as telekinetic force, localized environmental disturbances, and, in advanced cases, violent possession of the living.

\noindent\textbf{Stats}: Tier III‑IV, Cap 3–5. Incorporeal, Raging, Hungry.\\
\textbf{Harm}: 4. \textbf{Fatigue}: Spirit.

\noindent\textbf{Traits}: 
\emph{Telekinetic Throw} (DV 3, Harm 2), \emph{Possession Attempt} (DV 4, Spirit+Resist or Controlled).

\noindent\textbf{Complications}: Aura of Fear – all actions Desperate within Near range. Environmental Collapse – timers advance.

\noindent\textbf{Resolution}: Containment ritual (The Sentinel’s mirror rite) followed by negotiation with the underlying Anchor. The rite does not destroy the Poltergeist; it returns it to the state of an Anchor, at which point standard negotiation may resume.

\noindent\textbf{Clock}: \emph{Poltergeist Escalation [8]} – when filled, the spirit reaches Stage Four (The Mask): possession of a living host. Exorcism at this stage is dangerous and often fatal to the host.

\gray{I have seen a Poltergeist in Stage Four. It wore the face of a child. The child had been dead for three days. The host was the mother. Do not let it fill.}

\subsection*{Ghoul (Starved Anchor)}

\noindent\textbf{Description}: An Anchor that is denied resolution for too long may begin to consume the life force of the living to sustain its own existence. These ghouls retain their intelligence and memory. They are the most dangerous, because they are the most pitiable.

\noindent\textbf{Stats}: Tier II‑III, Cap 2–4. Corporeal, Hungry, Intelligent.\\
\textbf{Harm}: 4. \textbf{Fatigue}: Body.

\noindent\textbf{Traits}: 
\emph{Claw} (DV 2, Harm 1 + Fatigue 1), \emph{Bite} (DV 3, Harm 2 + poison), \emph{Regeneration} – ignore first 2 Harm each scene, \emph{Ravenous} – must feed or frenzy.

\noindent\textbf{Complications}: A ghoul that has not fed for a month will enter a state of torpor, appearing dead. It can be buried safely, but if disturbed, it will wake ravenous.

\noindent\textbf{Resolution}: Silver weapon to the heart, or fire. Negotiation possible if the ghoul is fed. There is no known cure for ghoulism; the Well of Ichor in Zahk’Kozahn offers release, but the price is memory.

\noindent\textbf{Clock}: \emph{Hunger [6]} – when filled, the ghoul attacks indiscriminately. The hunger is not a choice; it is a habit.

\saikou{I negotiated a feeding schedule with a ghoul in Port Shed. Animal blood, provided by a butcher, in exchange for her labor as a night watchman. She has not harmed a human in five years. Her mother knows. The butcher knows. I have told no one else until now.}

\subsection*{Spectral Procession (Class III)}

\noindent\textbf{Description}: A Procession is not a single spirit, but a route – a path that the dead walk, sometimes for centuries, between two points of significance. Battlefields, plague pits, shipwreck coasts. The spirits do not interact with observers. They are marching, not haunting. Do not stand in their way.

\noindent\textbf{Stats}: Tier III, Cap 4 per spirit (but treat as a swarm). Incorporeal, Processional, Unresponsive.\\
\textbf{Harm}: Cannot be fought; each spirit has Harm 2, but they move in a column of 1d20.

\noindent\textbf{Traits}: 
\emph{Marching} – the spirits follow a fixed route that predates current roads. \emph{Unresponsive} – they do not speak or react unless interrupted.

\noindent\textbf{Complications}: Interrupting a Procession – by standing in its path, attempting to speak to its members, or building over its route – provokes a violent response from the entire column. The response is not anger; it is the reflex of a machine.

\noindent\textbf{Resolution}: Observe, document, and – if the route must be altered – negotiate with the destination, not the dead. Moving the terminus (a grave, a shrine) will often reroute the spirits.

\noindent\textbf{Clock}: \emph{Procession’s Gait [8]} – ticks each time the route is blocked or diverted. When full, the column becomes a Poltergeist swarm.

\dusana{I once saw a Procession of drowned sailors walk through a market at midnight. They did not look at the stalls, the fish, the living. They looked only at the sea. The next day, the market smelled of salt for a week.}

---

\section*{New Variants from the Margins}

The following entries are not in Saikou Ira’s original taxonomy. I have added them from my own observations – and from the whispers I have heard in the little room with the salt‑ringed table. Treat them with caution. Treat them with respect. Treat them as you would a debt you cannot afford.

\subsection*{The Counting Wraith}

\noindent\textbf{Description}: Appears when someone miscounts – nine steps, nine cups, nine breaths. It is not a ghost of a person; it is a ghost of an \textbf{error}. The error is its anchor; the witness is its prey.

\noindent\textbf{Stats}: Tier III, Cap 3. Incorporeal, Patient, Hungry for Names.\\
\textbf{Harm}: 3. \textbf{Fatigue}: Spirit.

\noindent\textbf{Traits}: 
\emph{Name Theft} (DV 3, Fatigue 1) – the Wraith takes a single minor memory: the name of a street, the taste of bread, the face of a porter. \emph{Count Correction} – If the victim recites aloud the exact sequence of eight steps they took to cross the threshold where the Wraith appeared, the Wraith recedes. If they miscount, it returns the next night.

\noindent\textbf{Complications}: The Wraith cannot be fought; it can only be \textbf{counted}. A silver mirror will show its outline, but not its face.

\noindent\textbf{Resolution}: A correct count. If the count is wrong, the Wraith takes another memory. The memory is always small, necessary, and noticed only after it is gone.

\noindent\textbf{Clock}: \emph{Wraith’s Patience [6]} – ticks each time the victim miscounts. When full, the Wraith takes a \textbf{significant} memory: a name, a face, a year. I have seen a merchant forget his daughter’s name. He did not weep. He could not remember why he should.

\lerris{I lost the smell of rain on dry stone for a week after my first encounter. I did not miss it until it was gone. Now I count my breaths in eights, even when I am alone. The ninth breath is not for me.}

\subsection*{The Oath‑Keeper}

\noindent\textbf{Description}: A spirit of unpaid obligation. It appears when a promise has been broken and not acknowledged. It cannot be fought. It cannot be bargained with – except through the payment of the original debt.

\noindent\textbf{Stats}: Tier IV, Cap 5. Incorporeal, Patient, Inexorable.\\
\textbf{Harm}: Cannot be harmed. \textbf{Fatigue}: None.

\noindent\textbf{Traits}: 
\emph{Demand} – The Oath‑Keeper speaks a single sentence: “You owe.” It does not specify the debt; the debtor must remember. \emph{Interest} – Each time the Oath‑Keeper appears, the original debt grows. A forgotten promise to return a book becomes a curse on the bloodline.

\noindent\textbf{Complications}: The Oath‑Keeper will not harm the debtor directly. It will harm those the debtor loves. It is not cruel; it is \textbf{efficient}.

\noindent\textbf{Resolution}: Pay the original debt – in coin, in kind, or in a story that serves as interest. The payment must be witnessed. If the original debtor is dead, the debt passes to the heir. I have seen a son pay his father’s gambling debt to a ghost. The ghost accepted the coin and vanished. The son’s own debts were forgiven.

\noindent\textbf{Clock}: \emph{Debt’s Interest [8]} – ticks each time the debt is ignored. When full, the Oath‑Keeper claims the debtor’s name. The debtor becomes a stranger to everyone who knew them.

\saikou{The Oath‑Keeper is not a demon. It is a \textbf{creditor}. I have seen a village pay a collective debt with a single story told around the fire. The Walker listened, nodded, and did not return. The story was about a woman who forgot to name her child. The child’s name was the debt.}

\subsection*{The Weeping Shade (Class II – variant)}

\noindent\textbf{Description}: A ghost bound by a secret promise of silence. It cannot speak its unfinished business; it can only weep. The weeping is not sadness; it is the pressure of an unspoken name.

\noindent\textbf{Stats}: Tier II, Cap 2. Incorporeal, Bound, Mute.\\
\textbf{Harm}: 2. \textbf{Fatigue}: Spirit.

\noindent\textbf{Traits}: 
\emph{Weeping} – all within Near range suffer −1 die to Empathy and Insight rolls (the weeping is distracting). \emph{Touch of Silence} – on a successful Manifestation (DV 3), the Shade touches a target; the target cannot speak for one exchange.

\noindent\textbf{Complications}: The Shade cannot be reasoned with directly; it cannot speak. It can only be \textbf{named}.

\noindent\textbf{Resolution}: A public confession. Someone who knows the secret the Shade is bound to protect must speak it aloud in the presence of a witness. The Shade will hear, nod, and pass. I have seen a woman confess to a murder that was not her own; the Shade had been bound by the real murderer. The Shade passed; the woman became haunted by a different silence.

\noindent\textbf{Clock}: \emph{Silence’s Weight [6]} – ticks each time the secret is kept. When full, the Shade manifests permanently in the room where the promise was made. The room becomes a threshold; the threshold becomes a prison.

\dusana{The Weeping Shade in the hospice on Deadman’s Quay was a midwife who had promised not to reveal a stillbirth. I spoke the child’s name. The Shade stopped weeping. It did not thank me. It did not need to.}

---

\section*{Exorcism & Negotiation (Quick Reference)}

Not all unquiet dead require steel. Some require \textbf{witness}. Use this table as a guide.

| Class | Approach | Price | Failure Consequence |
|-------|----------|-------|---------------------|
| Echo | Speak the truth of the death aloud | None, but must be witnessed | Echo expands its range |
| Anchor | Fulfill unfinished business or substitution ritual | Memory, name, or object | Anchor becomes Poltergeist |
| Poltergeist | Sentinel’s mirror rite + negotiation with underlying Anchor | One memory per participant | Possession of a host |
| Ghoul | Silver or fire, or negotiated feeding schedule | Animal blood or a promise | Frenzy, indiscriminate attack |
| Procession | Redirect the terminus (grave, shrine) | A story or a coin | Column becomes Poltergeist swarm |
| Counting Wraith | Correct count of eight steps | Minor memory (if miscount) | Theft of significant memory |
| Oath‑Keeper | Pay original debt in coin, kind, or story | Witnessed payment | Debtor’s name is erased |
| Weeping Shade | Public confession of the secret | A name spoken aloud | Shade becomes permanent anchor |

\lerris{I keep this table pressed between two blank leaves. The leaves are not blank; they are waiting for the name I have not yet learned to speak. When I add it, I will be done.}

\gray{Exorcism is negotiation. Negotiation is memory. Memory is debt. Pay attention to the first line.}
\saikou{I have performed the Sentinel’s rite seven times. The mirrors are always warm afterward. The shards are always one fewer than I started with.}
\dusana{I have paid a ghost with a story. The story was about a fox who forgot its name. The ghost laughed and left. The fox is still out there, somewhere, waiting to be remembered.}

